\documentclass[10pt,a4paper]{amsart}
\usepackage[margin=19mm]{geometry}
\usepackage{amsmath,amssymb,mathtools}
\usepackage[scr=rsfs,cal=euler]{mathalfa}
\usepackage{xfrac}
\usepackage[unicode,hidelinks,bookmarks=false]{hyperref}
\newcommand{\ie}{i.e.\ }
\newcommand{\cf}{cf.\ }
\newcommand{\resp}{respectively\ }
\newcommand{\Subsection}{\subsection}
\providecommand{\nobreakdash}{\nobreak\hbox{-}}
\setlength{\emergencystretch}{2em}
\multlinegap0pt
\usepackage{bm}
\usepackage{bbm}
\usepackage{dsfont}
\usepackage{xspace}
\usepackage{cancel}
\usepackage{mathtools}
\usepackage{cleveref}
\usepackage{amsthm}
\makeatletter
\@ifundefined{theorem}{\newtheorem{theorem}{Theorem}}{}
\renewenvironment{proof}[1]{\par\medskip\noindent\textbf{#1}\enspace\ignorespaces}{\par\medskip}
\makeatother
\title[On Locational Marginal Emissions in Electricity Markets: A T]{On Locational Marginal Emissions in Electricity Markets: A Two-Layered Dispatch Mechanism and Its Fundamental Theorems}
\author{Luc Cote, Andy Sun}
\date{Source: arXiv:2603.19530v1 (CC BY 4.0)}
\begin{document}
\maketitle

\noindent\emph{Source and licence.} On Locational Marginal Emissions in Electricity Markets: A Two-Layered Dispatch Mechanism and Its Fundamental Theorems. Version arXiv:2603.19530v1 (CC BY 4.0), distributed under the Creative Commons Attribution 4.0 International licence, as recorded in the arXiv licence metadata for this exact version and shown on the abstract page https://arxiv.org/abs/2603.19530v1. Full rights notice: licenses/arxiv-2603.19530-CC-BY-4.0.txt.\par\smallskip

\noindent\textit{Editorial scope.} Counted unit: the Theorem whose source label is \texttt{thm:carbonfootprint} in the article (source0.tex lines 589--594, source label \texttt{thm:carbonfootprint}), delivered together with the article's own complete original proof of it (source0.tex lines 595--653, one \texttt{proof} environment of 59 lines). No mathematics is written, repaired or re-derived: statement and proof are the article's own text line for line. Required context retained uncounted: p:dcOPF (lines 187--196); p:OPF (lines 219--222); eq:cm:fb (lines 225--238); eq:cm:kvl (lines 225--238); eq:dual:f (lines 474--482); eq:dual:voltage (lines 474--482). Every definition, display and label of those lines is retained unchanged. Omitted from this package and not redistributed: 1--186, 197--218, 223--224, 239--473, 483--588, 654--2055. Nothing inside a retained excerpt is omitted or altered: every retained body is delivered exactly as the article's lines are written, so no omission receipt of any kind accompanies this package. Citations: the retained excerpts carry no citation command at all, so no bibliography entry of the article is transcribed and no reference is redistributed. Reference adaptation: none was needed and none was made; every reference inside the delivered unit resolves inside it, so no unresolved reference or citation is produced. Printed slips kept literal: no printed word, symbol, hypothesis or number of the counted statement or of its proof was altered. Layout and class: the article's own class is \texttt{article} with packages including geometry, natbib, inputenc, ifthen, float, graphicx, nomencl, authblk, hyperref, amsmath, amssymb, amsthm, cleveref, pgfplots; the wrapper is the lane's pinned \texttt{amsart} editorial wrapper at 10pt on A4, which restates the theorem-like environments the retained text needs and adds the article's own macro definitions that the retained text needs (none), with extra wrapper packages (bm, bbm, dsfont, xspace, cancel, mathtools, cleveref). The delivered numbering is the unit's own. Assets: no retained span contains a figure environment, a graphics-inclusion command, a PDF-page inclusion, a table or a TikZ picture; no image file is copied into this package, the package declares no asset dependency and no image file is redistributed with it. Licence: version arXiv:2603.19530v1 of this article is distributed under the Creative Commons Attribution 4.0 International licence, as recorded in the arXiv licence metadata for this exact version and shown on the abstract page https://arxiv.org/abs/2603.19530v1. A full rights notice is delivered at licenses/arxiv-2603.19530-CC-BY-4.0.txt. Characters: every retained excerpt is pure ASCII, so the delivered engine, XeLaTeX with its default Latin Modern text font, reports no missing character.
\begin{align}
   DCOPF(P^D):= \min_{P^G, \theta, f} \quad  & c^T P^G \tag{DCOPF Model} \label{p:dcOPF}\\
    \text{s.t.}\quad &
    B_{ij} (\theta_i - \theta_j) - f_\ell =  0, &\forall \ell=(i,j) \in L,&~[z_\ell],\label{eq:dcopf:angle}\\
   &  \sum_{g \in G_i} P^G_g - \sum_{\ell = (i,j) \in L}f_\ell + \sum_{\ell = (j,i) \in L}f_\ell - P_i^D= 0, & \forall i \in N,&~[\pi_i], \label{eq:dcopf:fb}\\
    & f_\ell \leq f^{max}_\ell, & \forall \ell \in L,&~[\rho^+_\ell],  \label{eq:dcopf:fmax}\\
    & f_\ell \geq -f^{max}_\ell, & \forall \ell \in L,&~[\rho^-_\ell],  \label{eq:dcopf:fmin}\\
    & P^G_g \leq P^{max}_g, & \forall g \in G,&~[\gamma^+_g],\label{eq:dcopf:gmax}\\
    & P^G_g \geq P^{min}_g, & \forall g \in G,&~[\gamma^-_g]\label{eq:dcopf:gmin}.
\end{align} 

\begin{align}
   E(P^D):= \min_{P^G, \theta, f}\quad  & \sigma^T P^G \tag{Combined Model} \label{p:OPF}\\
    \text{s.t.}\quad &   (P^G, f, \theta) \in X^*(P^D),
\end{align}

\begin{align}
    &B_{ij} (\theta_i - \theta_j) - f_\ell =  0, &\forall \ell=(i,j) \in L, &~[p_{z,\ell}],\label{eq:cm:kvl}\\
    &\sum_{g \in G_i} P^G_g - \sum_{\ell = (i,j) \in L}f_\ell + \sum_{\ell = (j,i) \in L}f_\ell - P_i^D= 0, & \forall i \in N, &~[p_{\pi,i}],\label{eq:cm:fb}\\
    &f_\ell \leq f^{max}_\ell, & \forall \ell \in L,&~[p^+_{\rho,\ell}], \label{eq:cm:fp}\\
    &f_\ell \geq -f^{max}_\ell & \forall \ell \in L,&~[p^-_{\rho,\ell}], \label{eq:cm:fm}\\
    &P^G_g \leq P^{max}_g, & \forall g \in G, &~[p^+_{\gamma,g}],\\
    &P^G_g \geq P^{min}_g, & \forall g \in G, &~[p^-_{\gamma,g}],\\
    &\gamma^+_g + \gamma^-_g + \sum_{i \mid g \in G_i} \pi_i  = c_g, &\forall g \in G, \\
    &-z_\ell + \rho^+_\ell + \rho^-_\ell - \pi_i + \pi_j = 0, &\forall \ell=(i,j)\in L, \\
    &\sum_{\ell = (i,j) \in L} B_\ell z_\ell - \sum_{\ell = (j,i) \in L} B_\ell z_\ell = 0, &\forall i \in N, \\
    &\rho^+, \gamma^+ \leq 0, \\
    &\rho^-, \gamma^- \geq 0, \\
    &c^TP^G - \pi^T P^D - (\rho^+ - \rho^-)^Tf^{max} - (\gamma^+)^T P^{max} - (\gamma^-)^T P^{min} \le 0, &&~[p_o]. \label{eq:cm:opt}
\end{align}

    \begin{align}
        &\hat{p}_{\pi,i}-\hat{p}_{\pi,j}-\hat{p}_{z,\ell} - \hat{p}^+_{\rho,\ell} - \hat{p}^-_{\rho,\ell} =0 ,&\forall \ell=(i,j)\in L, &~[f_\ell], \label{eq:dual:f}\\
        &\sum_{\ell = (i,j)}B_{\ell} \hat{p}_{z,\ell}-\sum_{\ell = (j,i)}B_{\ell}  \hat{p}_{z,\ell} =0 ,&\forall i \in N, &~[\theta_i], \label{eq:dual:voltage} \\
        &\sigma_g-\hat{p}_{\pi,g}-\hat{p}^+_{\gamma,g}-\hat{p}^-_{\gamma,g}-\hat{p}_o c_g=0, &\forall g \in G, &~[P^G_g]\label{eq:dual:pg}\\
        &\hat{p}^+_{\gamma}, 
        \hat{p}^+_{\rho} \leq 0, \label{eq:dual:plus}\\
        &\hat{p}^-_{\gamma}, 
        \hat{p}^-_{\rho} \geq 0.\label{eq:dual:minus}
    \end{align}

\begin{theorem}[Carbon Footprint Theorem] \label{thm:carbonfootprint}
The carbon accounts of all loads, generators, and transmission lines sum to the total system carbon emissions:
\begin{equation} \label{eq:carbonfootprint}
    \sum_{g \in G} \sigma_g P^G_g = \sum_{i \in N}LME_i P^D_i + \sum_{g \in G}(\sigma_g - LME_g) P^G_g + \sum_{\ell \in L} SCI_\ell |f_\ell|.
\end{equation}
\end{theorem}

\begin{proof}{of \Cref{thm:carbonfootprint}}
    We begin by defining an arbitrary set of optimal primal and dual solutions to the \ref{p:OPF} of the form $P^G$,$f$, $\theta$, $\pi$, $\gamma^+$, $\gamma^-$,$\rho^+$, $\rho^-$, $z$, and $p$.
    % \begin{equation}
    %     \text{maximize} \quad \sum_{\ell \in L} (p^+_{\rho,\ell}-p^-_{\rho,\ell})f^{max} + \sum_{g \in G} p^+_{\gamma,g} P^{max}_g + p^-_{\gamma,g}P^{min}_g - \sum p_{PG,g} c_g + \sum_{i \in N} p_{\pi,i}P^D_i \tag{Combined Dual} \label{p:CDualOPF}
    % \end{equation}
    % \begin{align}
    %     \text{s.t. }
    %     -p_{o} c_g + p_{\pi,g} + p^+_{\gamma,g} + p^-_{\gamma,g} &= \sigma_g, &\forall g \in G, \\
    %     B_{ij} (p_{\theta,i} - p_{\theta,j}) - p_{f,\ell} &=  0, &\forall \ell=(i,j) \in L,\\
    %     \sum_{g \in G_i} p_{PG,g} - \sum_{\ell = (i,j) \in L}p_{f,\ell} + \sum_{\ell = (j,i) \in L}p_{f,\ell} - p_o P^D_i&= 0, & \forall i \in N, \\
    %     p_{f,\ell} &\leq p_of^{max}_\ell, &\forall \ell \in L, \\
    %     p_{f,\ell} &\geq -p_of^{max}_\ell, &\forall \ell \in L,\\
    %     p_{PG,g} &\leq p_oP^{max}_g, &\forall g \in G,\\
    %     p_{PG,g} &\geq p_oP^{min}_g, &\forall g \in G,\\
    %     p_{z,\ell} + p^+_{\rho,\ell} + p^-_{\rho,\ell} - p_{\pi,i} + p_{\pi,j} &= 0, &\forall \ell=(i,j) \in L, \label{eq:cdual:fd} \\
    %     \sum_{\ell = (i,j) \in L} B_\ell p_{z,\ell} - \sum_{\ell = (j,i) \in L} B_\ell p_{z,\ell} &= 0, &\forall i \in N, \label{eq:cdual:td}\\
    %     p^+_{\rho}, p^+_{\gamma} &\leq 0, \\
    %     p^-_{\rho}, p^-_{\gamma} &\geq 0.
    % \end{align}
    We now proceed to evaluate the sum of carbon accounts:
    \begin{equation}
        \sum_{i \in N}LME_i P^D_i + \sum_{g \in G}(\sigma_g - LME_g) P^G_g + \sum_{\ell \in L} SCI_\ell |f_\ell|.
    \end{equation}
    Plugging in definition of $LME_i$ and $SCI_i$ gives
    \begin{equation}
        \sum_{i \in N}\left((p_{\pi,i}+p_o\pi_i) P^D_i\right) + \sum_{g \in G}\left((\sigma_g - p_{\pi,g}-p_o\pi_g) P^G_g\right) + \sum_{\ell \in L} \left( (p^+_{\rho,\ell}-p^-_{\rho,\ell}+p_o(\rho_\ell^+-\rho_\ell^-))|f_\ell|\right) 
    \end{equation}
    which from complementary slackness between $f_\ell$ and its corresponding dual variables for both the \ref{p:OPF} and the \ref{p:dcOPF} leads to
    \begin{equation}
        \sum_{i \in N}\left((p_{\pi,i}+p_o\pi_i) P^D_i\right) + \sum_{g \in G}\left((\sigma_g - p_{\pi,g}-p_o\pi_g) P^G_g\right) + \sum_{\ell \in L} \left( (p^+_{\rho,\ell}+p^-_{\rho,\ell}+p_o(\rho_\ell^++\rho_\ell^-))f_\ell \right) 
    \end{equation}
    and using strong duality of the \ref{p:dcOPF} simplifies to
    \begin{equation}
       \sum_{i \in N}\left(p_{\pi,i} P^D_i\right) + \sum_{g \in G}\left((\sigma_g - p_{\pi,g}) P^G_g\right) + \sum_{\ell \in L} \left( (p^+_{\rho,\ell}+p^-_{\rho,\ell})f_\ell \right). 
    \end{equation}
    Substituting $(p^+_{\rho,\ell}+p^-_{\rho,\ell})$ using \eqref{eq:dual:f} and rearranging leads to
    \begin{equation}
        \sum_{i \in N}\left(p_{\pi,i} (P^D_i-\sum_{g \in G_i}P^G_g+\sum_{\ell=(i,j)\in L}f_\ell-\sum_{\ell=(j,i)\in L}f_\ell)\right) + \sum_{g \in G}\left(\sigma_g P^G_g\right) - \sum_{\ell=(i,j) \in L} \left(p_{z,\ell}f_\ell \right).
    \end{equation}
    Where plugging in \eqref{eq:cm:kvl} and \eqref{eq:cm:fb} simplifies to
    \begin{equation}
        % \sum_{g \in G}\left(\sigma_g P^G_g\right) - \sum_{\ell=(i,j) \in L} \left(p_{z,\ell}f_\ell \right).
        \sum_{g \in G}\left(\sigma_g P^G_g\right) - \sum_{i\in N}\theta_i \left(\sum_{\ell=(i,j)\in L}p_{z,\ell} B_{\ell}-\sum_{\ell=(j,i)\in L}p_{z,\ell}B_{\ell} \right)
    \end{equation}
    % Finally, plugging in \eqref{eq:cm:kvl} rearranging the right sum gives
    % \begin{equation}
    %     \sum_{g \in G}\left(\sigma_g P^G_g\right) - \sum_{i\in N} \left(\sum_{\ell=(i,j)\in L}p_{z,\ell} B_{\ell}-\sum_{\ell=(j,i)\in L}p_{z,\ell}B_{\ell} \right)
    % \end{equation}
    which using \eqref{eq:dual:voltage} simplifies to
    \begin{equation}
        \sum_{g \in G}\sigma_g P^G_g.
    \end{equation}
    % substituting $\sigma_g$ using the first constraint of \Cref{p:CDualOPF} and rearranging leads to 
    % \[\sum_{i \in N}\left(p_{\pi,i} P^D_i\right) + \sum_{g \in G}\left((p^+_{\gamma,g}+p^-_{\gamma,g}-p_oc_g) P^G_g\right) + \sum_{\ell \in L} \left((p^+_{\rho,\ell}-p^-_{\rho,\ell})f_\ell\right)\]
    % finally, strict complementarity allows us to replace $P^G_g$ and $f_\ell$ with the appropriate limits thus yielding
    % \[\sum_{i \in N}\left(p_{\pi,i} P^D_i\right) + \sum_{g \in G}\left(p^+_{\gamma,g}P^{max}_g+p^-_{\gamma,g}P^{min}_g -p_oc_gP^G_g\right) + \sum_{\ell \in L} \left((p^+_{\rho,\ell}-p^-_{\rho,\ell})f^{max}_\ell\right)\]
    % which from LP strong duality is precisely equal to 
    % \[\sum_{g \in G} \sigma_g P^G_g\]
\end{proof}

\end{document}
